%% kapa: exam=2 round=30 year=2019 period=2 subject=appraisal_theory subject_ko=감정평가이론 seq=2 number=2 points=30
%% source: Q-Net 감정평가사 기출문제 (cst00309) artlSeq=5205595 file=raw/1008/5205595/감정평가이론.pdf
%% conversion: kordoc 4.24.51 -> tools/gen2.py (auto); review=proofread-gpt-terra; flags=none
\begin{longquestion}{【문제 2】}
시장가치에 대하여 다음의 물음에 답하시오. (30점) \par
\medskip\noindent 물음 1) ‘성립될 가능성이 가장 많은 가격(the most probable price)’이라는 시장가치의 정의가 있다. 이에 대해 설명하시오. (10점) \par
\medskip\noindent 물음 2) 부동산거래에 있어 ‘최고가격(highest price)’과 ‘성립될 가능성이 가장 많은 가격’을 비교ㆍ설명하시오. (10점) \par
\medskip\noindent 물음 3) 가치이론과 가치추계이론의 관계에 대해 각 학파의 주장내용과 이에 관련된 감정평가방법별 특징을 설명하시오. (10점) \par
\end{longquestion}
